\documentclass[oneside]{amsart}
\usepackage[all]{xy}
\usepackage{tikz-cd}
\usepackage[T1]{fontenc}
\usepackage{xstring}
\usepackage{xparse}
\usepackage{xr-hyper}
\usepackage{xcolor}
\definecolor{brightmaroon}{rgb}{0.76, 0.13, 0.28}
\usepackage[linktocpage=true,colorlinks=true,hyperindex,citecolor=blue,linkcolor=brightmaroon]{hyperref}
\usepackage[nameinlink]{cleveref}
\usepackage[left=1.25in,right=1.25in,top=0.75in,bottom=0.75in]{geometry}
\setlength{\marginparwidth}{2cm}
%\usepackage[charter,greekfamily=didot]{mathdesign}
%\usepackage{Baskervaldx}
\usepackage{amssymb}
\usepackage{stmaryrd}
\usepackage{mathrsfs}
\usepackage{mathpazo}
\usepackage{eulervm}
\linespread{1.05}

\usepackage[nobottomtitles]{titlesec}
\usepackage{marginnote}
\usepackage{enumerate}
\usepackage{longtable}
\usepackage{aurical}
\usepackage{microtype}
\usepackage{mathtools}

\newtheoremstyle{ega-env-style}%
{}{}{\rmfamily}{}{\bfseries}{.}{ }{\thmnote{(#3)}}%

\newtheoremstyle{ega-thm-env-style}%
{}{}{\itshape}{}{\bfseries}{. --- }{ }{\thmname{#1}\thmnumber{ #2}\thmnote{ (#3)}}%

\newtheoremstyle{ega-defn-env-style}%
{}{}{\rmfamily}{}{\bfseries}{. --- }{ }{\thmname{#1}\thmnumber{ #2}\thmnote{ (#3)}}%

\theoremstyle{ega-env-style}
\newtheorem*{env}{---}

\theoremstyle{ega-thm-env-style}
\newtheorem{theorem}{Theorem}[section]
\newtheorem{proposition}[theorem]{Proposition}
\newtheorem{lemma}[theorem]{Lemma}
\newtheorem{corollary}[theorem]{Corollary}
\newtheorem{conjecture}[theorem]{Conjecture}

\theoremstyle{ega-defn-env-style}
\newtheorem{definition}[theorem]{Definition}
\newtheorem{example}[theorem]{Example}
\newtheorem*{examples}{Examples}
\newtheorem*{remark}{Remark}
\newtheorem*{remarks}{Remarks}
\newtheorem*{notation}{Notation}
\newtheorem*{exercise}{Exercise}
\newtheorem*{properties}{Properties}
\newtheorem*{consequences}{Consequences}
\newtheorem*{note}{Note}
\newtheorem*{fact}{Fact}
\newtheorem*{claim}{Claim}

\makeatletter
\def\l@subsection{\@tocline{2}{0pt}{2.5pc}{2.2pc}{}}
\def
\section{\@startsection{section}{1}%
	\z@{.7\linespacing\@plus\linespacing}{.5\linespacing}%
{\normalfont\bfseries\Large\scshape\centering}}
\renewcommand{\@seccntformat}[1]{%
	\ifnum\pdfstrcmp{#1}{section}=0\textsection\fi%
\csname the#1\endcsname.~}
\makeatother

\def\mathcal{\mathscr}
%% Fonts

\def\sh{\mathcal}                   % sheaf font
\def\bb{\mathbf}                    % bold font
\def\cat{\mathsf}                   % category font
\def\numfont{\mathbb}

%% Font Letters

\def\CA{\mathcal{A}}
\def\CE{\mathcal{E}}
\def\CF{\mathcal{F}}
\def\CG{\mathcal{G}}
\def\CL{\mathcal{L}}
\def\OO{\mathcal{O}}
\def\CX{\mathcal{X}}
\def\b1{\mathbb{1}}

%% Cohomology

\def\CH{\mathrm{H}}                 % cohomology H
\def\CHH{\check{\HH}}               % Čech cohomology H
\def\RD{\mathrm{R}}                 % right derived R
\def\LD{\mathrm{L}}                 % left derived L
\def\dual#1{{#1}^\vee}              % dual
\def\Tor{\operatorname{Tor}}        % Tor
\def\Ext{\operatorname{Ext}}        % Ext
\def\HdR{\mathrm{H}_{\mathrm{dR}}}  % de Rham cohomology
\def\Zc{\underline{\mathbb{Z}}}     % constant sheaf with integer coeffs

%% Categories

\def\A{\cat{A}}                     % category A, usually abelian
\def\C{\cat{C}}                     % category C
\def\op{^\cat{op}}                  % opposite category
\def\Set{\cat{Sets}}                % category of sets
\def\Grp{\cat{Gps}}                 % category of groups
\def\Alg{\cat{Alg}}                 % category of algebras
\def\QCoh{\cat{QCoh}}               % category of quasicoherent sheaves
\def\supertilde{{\,\widetilde{\,}\,}}   % use \supertilde instead of ^\sim
\def\Map{\operatorname{Map}}        % morphisms (any category)
\def\Hom{\operatorname{Hom}}        % morphisms (additive category)
\def\iHom{\underline{\Hom}}         % interal hom
\def\End{\operatorname{End}}        % endomorphisms
\def\Aut{\operatorname{Aut}}        % automorphisms
\def\ker{\operatorname{ker}}        % kernel
\def\img{\operatorname{im}}         % image
\def\coker{\operatorname{coker}}    % cokernel
\def\pr{\operatorname{pr}}          % projection
\def\EssIm{\operatorname{EssIm}}    % essential image
\DeclareMathOperator*{\colim}{colim}   % colimit

\def\DAet{\cat{DA}\et}
\def\SH{\cat{SH}}
\def\Sh{\cat{Sh}}                   % category of sheaves
\def\PSh{\cat{PSh}}                 % category of presheaves

%% Schemes

\def\Proj{\operatorname{Proj}}      % Proj
\def\Supp{\operatorname{Supp}}      % support
\def\Spec{\operatorname{Spec}}      % Spec
\def\Spf{\operatorname{Spf}}        % formal Spec
\def\Aff{\mathbb A}                 % affine space
\def\P{\mathbb{P}}                  % projective space
\def\Pic{\operatorname{Pic}}        % Picard group

%% Standard Operators

\def\codim{\operatorname{codim}}    % codimension
\def\id{\operatorname{id}}          % identity

%% Arrows
\renewcommand{\to}{\mathchoice{\longrightarrow}{\rightarrow}{\rightarrow}{\rightarrow}}
\newcommand{\from}{\mathchoice{\longleftarrow}{\leftarrow}{\leftarrow}{\leftarrow}}
\let\mapstoo\mapsto
\renewcommand{\mapsto}{\mathchoice{\longmapsto}{\mapstoo}{\mapstoo}{\mapstoo}}
\def\isoto{\simeq}                  % isomorphism
\def\simto{\xrightarrow{\sim}}      % isomorphism arrow
\def\surjto{\twoheadrightarrow}     % surjetion
\def\injto{\hookrightarrow}         % injection

%% Under/Over Accents

\newcommand{\wh}[1]{\widehat{#1}}   % hat
\newcommand{\wt}[1]{\widetilde{#1}}    % tilde
\def\ul{\underline}                % underline

%% Spaces

\def\Z{\numfont Z}                  % integers
\def\Q{\numfont Q}                  % rationals
\def\N{\numfont N}                  % naturals
\def\R{\numfont R}                  % reals
\def\H{\mathcal H}

%% Groups

\def\GL{\bb{GL}}                   % general linear group
\def\SL{\bb{SL}}                   % special linear group
\def\Sp{\bb{Sp}}
\def\PSp{\bb{PSp}}
\def\GSp{\bb{GSp}}
\def\det{\operatorname{det}}       % determinant

%% Sub/Superscripts

\def\et{^\text{\'et}}              % \'etale
\def\an{^\text{an}}                % analytic
\def\th{^\text{th}}                % th

%% Misc

\newcommand{\defn}[1]{\textbf{#1}}  % definition highlighting
\def\defeq{:=}                     % definition equation
\def\eps{\varepsilon}              % correct epsilon
\newcommand{\gen}[1]{\left\langle\!\left\langle #1 \right\rangle\!\right\rangle}


\def\shHom{\sh{H}\textup{\kern-2.2pt{\Fontauri\slshape om}}}   % sheaf Hom
\def\shProj{\sh{P}\textup{\kern-2.2pt{\Fontauri\slshape roj}}} % sheaf Proj
\def\shExt{\sh{E}\textup{\kern-2.2pt{\Fontauri\slshape xt}}}   % sheaf Ext
\def\shGr{\sh{G}\textup{\kern-2.2pt{\Fontauri\slshape r}}}     % sheaf Gr
\def\shDer{\sh{D}\,\textup{\kern-2.2pt{\Fontauri\slshape er}}} % sheaf Der
\def\shDiff{\sh{D}\,\textup{\kern-2.2pt{\Fontauri\slshape if{}f}}\,} % sheaf Diff
\def\shHomcont{\sh{H}\textup{\kern-2.2pt{\Fontauri\slshape om.\,cont}}}   % sheaf Hom.cont
\def\shAut{\sh{A}\textup{\kern-2.2pt{\Fontauri\slshape ut}}}   % sheaf Aut
\def\shSym{\sh{S}\textup{\kern-2.2pt{\Fontauri\slshape ym}}}   % sheaf Sym

\makeatletter
\newcommand{\cbigoplus}{\DOTSB\cbigoplus@\slimits@}
\newcommand{\cbigoplus@}{\mathop{\widehat{\bigoplus}}}
\makeatother

\newcommand{\muT}{\mu_{\mathrm{Tam}}}
\newcommand{\F}{\mathbb{F}}
\DeclareMathOperator{\Bun}{Bun}
\DeclareMathOperator{\Tr}{Tr}
\DeclareMathOperator{\Frob}{Frob}
\DeclareMathOperator{\Ran}{Ran}
\newcommand{\red}{_{\mathrm{red}}}
\newcommand{\km}{\mathfrak{m}}
\DeclareMathOperator{\Sym}{Sym}
\DeclareMathOperator{\Shv}{Shv}
\DeclareMathOperator{\Gr}{Gr}
\def\D{\mathbb{D}}
\newcommand{\Sch}{\cat{Sch}}
\newcommand{\Fins}{\cat{Fin}^{\mathrm{s}}}
\newcommand{\Fun}{\cat{Fun}}
\DeclareMathOperator{\ins}{ins}

\mathchardef\mhyphen="2D
\newcommand{\iCat}{\cat{\infty\mhyphen Cat}}
\newcommand{\iCatSt}{\cat{\infty\mhyphen Cat}^{\mathrm{St}}}
\newcommand{\lMod}{\Lambda\cat{\mhyphen Mod}}

\def\surjfrom{\twoheadleftarrow}


\newcommand{\comp}{_{\mathrm{comp}}}	% completion sub
\newcommand{\D}{\mathrm{D}}
\newcommand{\prestk}{\mathrm{PreStk}}
\newcommand{\laxprestk}{\mathrm{LaxPreStk} }
\newcommand{\Cat}{\mathrm{Cat}}
\newcommand{\Ani}{\mathrm{Ani}}
\newcommand{\Fin}{\mathrm{Fin}}
\newcommand{\Sets}{\mathrm{Sets}}
\newcommand{\diag}{\mathrm{diag}}

\newtheorem{construction}[theorem]{Construction}
\usepackage[backend=biber,style=alphabetic,sorting=ynt]{biblatex}
\addbibresource{pwg.bib}

\title{The unital Ran spaces}
\author{Runlei Xiao}
\date{Padova, March 31, 2026}
\begin{document}
\maketitle

\setcounter{section}{0}
\section{Conventions and Notations}
\begin{itemize}
    \item $k$ is an algebraic closed field.
    \item $\Sch$ is the category of schemes of finite type over $k$.
    \item $X$ is a separated scheme of finite type over $k$.
\end{itemize}

\section{The motivation to construct unital Ran spaces}

In the proof of the Atiyah-Bott formula from \cite{abott}, Gaitsgory obtains the cohomological product formula via
\begin{itemize}
    \item Non-abelian Poincaré duality;
    \item $C^*(\Ran, \mathcal A)^\vee \simeq C^*(\Ran, \mathcal{B})$
\end{itemize}
Today, we  focus on the general recipe to obtain the second isomorphism.

We now explain to you why you have to take this recipe, "Taking out of unit", to construct this isomorphism. From Micha\l's talk \cite[Theorem~1.3]{}\space, you might expect $\mathcal A$ satisfying cohomological estimate and $\mathbb{D}  \mathcal A \simeq \mathcal B$ and \[    
C^*_C(\Ran, \mathbb{D} \mathcal{A}) \simeq C^*_c(\Ran, \mathcal{B})
\]
But 
\begin{proposition}
    $\mathcal{A}$ doesn't satisfy cohomological estimate and $\mathbb{D} \mathcal A$ is zero.
\end{proposition}

So we can't directly apply the Verdier duality argument; we have to do a "take unit out" modification on $\mathcal A$. The idea of these modification come from the analogies between $\Shv(\Ran)$ with factorisation structure and the category of non-unital associative $dg-$algebras. We resume some facts of non-unital associative algebras and their variations first:

For the category $dg\Alg_k$ of associative dg-algebras (not necessary with unit), we have a contravariant self functor, Koszul duality\(
       KD\colon dg\Alg^{}_k \to dg\Alg^{}_k\)
        \[
            KD(A) \coloneqq \ker(Hom_{A^+}(k,k) \to Hom_k(k,k))
        \]  
        where $A^+$ denotes the unital algebra of $A \oplus k$. 
There are two important examples to explain why we need to take the unit out:
\begin{itemize}
    \item If $A$ is a unital $dg-$algebra forgetting its unit, $k$ is a direct summand of $A^+$. Therefore $KD(A) = 0$.
    \item Let $A$ be an augmented unital algebra $k[\epsilon]/\epsilon^2$ where $\epsilon$ is degree $-1$ element. Let $A^{\mathrm{red}}$ denote the augmented ideal $(\epsilon)$ of $A$. $KD(A) \simeq uk[[u]]$ and $HH_*(A^{\mathrm{red}})^\vee \simeq HH_*(uk[[u]])$
\end{itemize}

Keep this intuition in our mind, we denote 
\begin{notation}
    $\Ran_{untl}$ for the unital Ran spaces, the sheaves on them are analogs of unital $dg-$algebras,
    $\Ran_{untl,aug}$ for the unital augmented Ran spaces, the sheaves on them are analogs of augmented unital $dg-$algebras.
\end{notation}

we denote it by $\mathcal A^{\mathrm{red}}$ because: 
\begin{itemize}
    \item $\mathcal A^{\mathrm{\mathrm{red}}}$ satisfies cohomological estimate, the $!-$fiber of $\mathcal A^{\mathrm{red}}$ over the distinct points $\{ x_1, \dots ,x_n\} \in \Ran$ is \[  
        \otimes_{i=1,\dots, n} C^*_{\mathrm{red}}(\Gr_{x_i})
    \]
    \item $\mathcal B^{\mathrm{red}} \simeq \mathbb{D} A^{\mathrm{red}}$, the $!-$fiber of $\mathcal B^{\mathrm{red}}$ over the distinct points $\{ x_1, \dots ,x_n\} \in \Ran$ is\[   
    \otimes_{i= 1, \dots, n} C_{\mathrm{red}}^*(BG_{x_i})
    \]
\end{itemize}

\subsection{Limit and colimit}
%If my intuition is correct, this construction works for an enriched category also.
In the following, we say $\infty-$category for $(\infty,1)-$category.
We review the model of colimit and limit in the $\infty-$category. It will help us to imagine  "spaces" that are constructed later. We also disregard the size issue.

\begin{construction}
    Let $\mathrm{C}, \mathrm{D}$ be $\infty-$categories, $F\colon \C \to \D$. Applying the strightening/Unstrightening theorem, we have a cocartesian fibration $$U \colon \int_C F \to \C$$
\end{construction}

Replace $\D$ with the category $\Cat_\infty$ of $\infty-$categories, and there is the following characterization of the limit of $F$.



\begin{proposition}\cite[\href{https://kerodon.net/tag/05RX}{Tag 05RX}]{kerodon}
    The $\infty-$category of coCartesian sections $$\Fun^{CCart}_{/C}(C , \int_C F)$$ is limit of $F$
\end{proposition}

\begin{corollary}\cite[\href{https://kerodon.net/tag/05RY}{Tag 05RY}]{kerodon}
    If $F\colon \C \to \Cat_\infty$ factor through $\Ani$, $\lim F$ is an anima
\end{corollary}

\begin{proof}
    Because $\Ani \to \Cat_\infty$preserves small limit. The conceptual explanation is following:
    \[  
        \Fun^{CCart}_{/\C}(\C , \int_F\C) = \Fun_{/\C}(C, \int_F \C)
    \]
  All arrows are $U-$cocartesian edges, especially the non-isomorphism arrow in $\int_F \C$ is from $\C$, therefore also a $U-$coCartesian edge. Since $\int_C F \to C$ is a left fibration. Applying \cite[\href{https://kerodon.net/tag/01GY}{Tag 01GY}]{kerodon}, $\$Fun_{/\C}(C, \int_F \C)$ is a left fibration over $\Delta^0$, therefore a Kan complex \cite[\href{https://kerodon.net/tag/019D}{Tag 019D}]{kerodon}
\end{proof}

\begin{proposition}
    Let $W$ be the collection of all $U-$coCartesian edges in $\int_CF$, the $\infty-$category of \(
        \int_C F[W^{-1}]
    \) is the colimit of $F$
\end{proposition}

\begin{corollary}
      If $F\colon \C \to \Cat_\infty$ factor through $\Ani$, $\colim F$ is an anima.
\end{corollary}

\begin{proof}
    Because all edges in $\int_F C$ are coCartesian edges.
\end{proof}

\section{Lax prestacks}
%In the following, we consider $\Cat_\infty$ as an $(\infty)-$category.
\begin{definition}\cite[Paragraph~2.1]{abott}
    \emph{A lax prestack} is a functor \[   
    \mathcal{Y} \colon  \Sch_S^{op} \to \Cat_\infty.
    \] Applying the Straightening/Unstraightening theorem,$\mathcal{Y}$ induces Cartesian fibration \[ 
    \int_\Sch \mathcal{Y} \to \Sch
    \]
    by passing to fiberwise opposite categories.
\end{definition}

\begin{remark}
    If $\mathcal{Y}$ is a prestack, we have $\mathcal{Y} \simeq \mathcal{Y}^{op}$, because $\mathcal{Y}(S) \simeq \mathcal{Y}(S)^{op}$ for any Anima.
\end{remark}

\begin{proposition}
The fully faithful embedding of $\prestk$ in $\laxprestk$ admits a left-adjoint    \[
    \prestk \to \laxprestk
\] induced by the group completion\[    
    \Cat_\infty \to \Ani \space: \space \C \mapsto \C[\mathrm{Mor}(C)^{-1}]
\]
\end{proposition}

\subsection{Sheave on lax prestacks}
Let we write the "limit" description of sheaves on prestacks in Michal's more formal, namely

\begin{construction}
    Consider $\Shv^!\colon \Sch^{op} \to \Cat_\infty$, we obtain a Cartesian fibration \[   
    \int_{\Sch}  \Shv^! \to \Sch
    \]
\end{construction}

\begin{definition}
    For $\mathcal{Y} \in \laxprestk$, we denote $\Shv^!(\mathcal{Y})$ the category of horizontal  functors which make the following diagram commute.% https://q.uiver.app/#q=WzAsMyxbMCwwLCJcXGludF9cXG1hdGhjYWx7WX0gXFxTY2giXSxbMiwwLCJcXGludF97XFxTaHZeIX0gXFxTY2giXSxbMSwyLCJcXFNjaCJdLFswLDFdLFswLDJdLFsxLDJdXQ==
\[\begin{tikzcd}
	{\int_\Sch \mathcal{Y}} && {\int_{\Sch} \Shv^!} \\
	\\
	& \Sch
	\arrow[from=1-1, to=1-3]
	\arrow[from=1-1, to=3-2]
	\arrow[from=1-3, to=3-2]
\end{tikzcd}\]
i.e, an object $F \in \Shv^!(\mathcal{Y})$
\begin{itemize}
    \item  $S \in \Sch$, $y \in \mathcal{Y}(S)$, it associates to $F_{S,y} \in \Shv^!(S)$
    \item[New] $(y_1 \to y_2) \in \mathcal{Y}(S)$, it associates to $(F_{S,y_1} \to F_{S,y_2}) \in \Shv^!(S)$
    \item $f \colon S' \to S, y \in \mathcal{Y(S)}$, it associates to $F_{S',f(y)} \simeq f^!(F_{S,y}) \in \Shv^!(S)$
\end{itemize}
Therefore, we have $Shv^! \colon \laxprestk \to \Cat_\infty$
\end{definition}

\begin{proposition}
    For a morphism $f \colon \mathcal{Y}_1 \to \mathcal{Y}_2$, we have $f^!$ commute with colimit and limit.
\end{proposition}

\begin{proof}
For a sheaf $F$ on a lax prestack, it is determined by its pointwise value. And  $f^!$ is right adjoint to $f_!$,therefore $f^!$ preserves limit.

Let $F \in \Shv^!(\mathcal{Y}_2)$ we have $$(f^!\colim F^i)_{(S,y)} \simeq f^!(\colim F^i_{(S,y)}) \simeq \colim f^!F^i_{(S,y)}$$ over $\Shv^!(S)$, the last equivalence is because $f_!$ commute with compact objects and $Shv^!(S)$ is compactly generated.
\end{proof}

\section{Unital Ran space}
We will construct $\Ran_{untl}$ as a lax prestack.
\begin{notation}
    Let $\Fin$ denote the category of finite non-empty sets and all maps. For a set $A$, the functor
    \[  
    I \to A^I, \thickspace \Fin^{op} \to \Sets 
    \]
    defines a Cartesian fibration in groupoids $A^{\Fin} \to \Fin$
\end{notation}
We first construct a bigger lax prestack $X^{\Fin}$, then we obtain $\Ran_{untl}$ by localisation.
\begin{definition}
    Let $X^{\Fin}$ be the lax prestack that assigns to $S \in \Sch$ the category \[   
    \Map(S,X)^{\Fin}.
    \]
    We obtain a Cartesian fibration $\int_{\Sch \times \Fin} X^{\Fin} \to \Sch \times \Fin$ in groupoid and a Cartesian fibration $\int_{\Sch} X^{\Fin}  \to \Sch$.
\end{definition}

\begin{notation}
    For a set $A$, let $A^{\Fin,\sim}$ denote the localization of $A^{\Fin}$ by inverting all Cartesian arrows that lie above those surjection maps in $\Fin$. I,e, we identify  $f \colon I \to A $ by their image, and keep track of the conclusion relation.
\end{notation}

\begin{definition}
    Let $\Ran_{untl}$ denote the lax prestack $X^{\Fin, \sim}$ where $X^{\Fin, \sim} (S) = \Map(S,X)^{\Fin,\sim}$
\end{definition}

\begin{remark}
    In the paper \cite{abott}, a unital ran space $\Ran_{untl}$ is a functor associated to $S \in \Sch$ to $Subset_f(\Map(S,X))$ where $Subset_f$ denote for the category of finite subsets with embedding. These two definitions give the same description in one space, namely $F \simeq X^{Fin,\sim}$. To check this fact, we only need to prove $Subset_f\Map(S, X) \simeq \Map(S,X)^{\Fin,\simeq}$

    Since $\Map(S, X)^{\Fin} \to Subset_f(\Map(S,X))$ is a Cartesian fibration, and the cartesian arrow in $\Map(S,X)^{\Fin}$ over the surjective morphism in $\Fin$ get sent to the isomorphisms in $Subset_f(\Map(S,X))$. Moreover, for a subset $B \subset \Map(S,X)$, its fiber is  the slice category over $B$ generated by the  surjection $J \to B$, which has an initial object $\id: B \to B$, so the colimit of a category with an initial object is a point.
\end{remark}

\begin{definition}
    By construction $\Ran_{untl} \coloneqq X^{\Fin,\sim}$,  $\Shv^!(\Ran_{untl})$ is the full subcategory of $\Shv^!(X^{\Fin})$. So $F \in \Shv(\Ran_{untl})$ is an assignment
    \[  
        (J \in \Fin) \mapsto F_J \in \Shv(X^J)
    \]
    equipped with a homotopy-compatible system of morphisms
    \[  
        (I \xrightarrow[]{\alpha}J) \mapsto (\diag_\alpha)^!F_I \to F_J
    \]
    Moreover if $\alpha$ is a surjection in $\Fin$, $(diag_\alpha)^!F_I \to F_J$ is an isomorphism.
\end{definition}

\begin{remark}
    You will see that sheaves on $\Ran_{untl}$ encode more information from the $\alpha\colon I \to J$ injection.
\end{remark}

\section{Adjunction of forgetting and adding the unit}
We would like to construct $\mathrm{AddUnit}\colon \Shv(\Ran) \to \Shv(\Ran_{untl})$ such that for the distinct point  $I=\{x_1, \dots, x_n\}$, $\mathcal{A} \in \Shv(\Ran)$ \[    
    \mathrm{AddUnit}(A_{I}) = \oplus_{J \subset I} \mathcal A_J
\]

The construction of the functor $\mathrm{AddUnit}$ arise a lax prestack  $\Ran^{\to}$
\begin{definition}
    $\Ran^{\to}$ is a lax prestack. For each scheme $S \in \Sch$, $\Ran^{\to}(S)$ be the category 
    \begin{itemize}
        \item Objects are pairs $(J \subseteq I)$ of finite subsets of $\Map(S,X)$ with $J \neq \emptyset$
        \item Morphisms from $(J \subseteq I)$ to $(J_1 \subseteq I_1)$ are inclutions $I \subseteq I_1$ such that $J =J_1$
    \end{itemize}
\end{definition}

\begin{remark}
    Categorically, $\Ran^{\to} \simeq \Ran \times^{\to}_{\Ran_{untl}} \Ran_{untl}$ fiberwisely.
\end{remark}

We have the maps 
\[  
    \psi \colon \Ran^{\to} \to \Ran_{untl}, \thickspace (J\subseteq I) \mapsto I,
\]
\[  
    \xi \colon \Ran^{to} \to \Ran, \thickspace  (J \subseteq I) \mapsto J
\]
\[  
    \nu \colon \Ran \to \Ran^{\to}, \thickspace J \mapsto (J \subseteq J)
\]

\begin{proposition}\cite[Lemma~4.3.2]{abott}
     The map $\psi$ is pseudo-proper
\end{proposition}

\begin{proof}
    We first show $\Ran^{\to}(S) \to \Ran_{untl} (S)$ is a coCartesian fibration in groupoid for $S \in \Sch$. Let $I$ be an object of $\Ran_{untl}(S)$, the fiber of $\Ran^{\to}(S)$ over it is the groupoid of non empty subsets $J \subseteq I$. For a morphism $I_1 \subseteq I_2$ in $\Ran_{untl}(S)$, we have map between the corresponding fibers, given by 
    \[
    (J \subseteq I_1) \mapsto (J\subset I_2)
    \]

    So now we just check pseudo-properness fiberwisely. Fix $I \subset \Map(S,X)$, 
    \[  
    S \times_{\Ran_{untl}} \Ran^{\to} \to S
    \]
    is pseudo-proper, because for the index category over $I$, whose objects are the data of $J \to J' \twoheadleftarrow I$, where $J$ is a finite non-empty set. Morphisms are commutative diagrams that induce a surjection on $J$. Now the fiber product $  S \times_{\Ran_{untl}}\Ran^{\to} $ is
    \[  
        \colim_{J \to J' \twoheadleftarrow I} X^{J'} \times_{X^I} S.
    \]    
\end{proof}

\begin{definition}
    The functor 
    \[
    \mathrm{AddUnit} \colon \Shv(\Ran) \to \Shv(\Ran_{untl})
    \]
    to be the composition of $\psi_! \circ \xi^!$
\end{definition}

Explicitly, let $S \in \Sch$, let $I \subset \Ran_{untl}(S)$ corresponding to
\[  
    S \to X^{\circ, I} \to X^I \to \Ran \xrightarrow[]{\phi} \Ran_{untl}.
\]

\begin{proposition}
    For $F \in \Shv(\Ran)$, the corresponding object $\mathrm{AddUnit}(F)_{S,I} \in \Shv(S)$ is given by 
    \[  
        \oplus_{\emptyset \neq J \subseteq I} F_{S,J}
    \]
\end{proposition}

\begin{proof}
Consider a Cartesian diagram with a morphism $\xi$, apply \cite[Corollary~2.3.5]{abott} on $I^! \psi_! \xi^!$, we hold this result.
% https://q.uiver.app/#q=WzAsNSxbMiwwLCJcXFJhbl57XFx0b30iXSxbMiwyLCJcXFJhbl97dW50bH0iXSxbNCwyLCJTIl0sWzQsMCwiXFxzcWN1cF97SiBcXHN1YnNldGVxIEl9IFhfSiJdLFswLDIsIlxcUmFuIl0sWzIsMSwiSSIsMl0sWzAsMSwiXFxwc2kiXSxbMywwXSxbMywyXSxbMCw0LCJcXHhpIiwyXV0=
\[\begin{tikzcd}
	&& {\Ran^{\to}} && {\sqcup_{J \subseteq I} X_J} \\
	\\
	\Ran && {\Ran_{untl}} && S
	\arrow["\xi"', from=1-3, to=3-1]
	\arrow["\psi", from=1-3, to=3-3]
	\arrow[from=1-5, to=1-3]
	\arrow[from=1-5, to=3-5]
	\arrow["I"', from=3-5, to=3-3]
\end{tikzcd}\]

\end{proof}
Now we want to construct the forgetting functor by the following observation.

Consider the morphism 
\[
    \phi \colon \Ran \to \Ran_{untl}
\]
by the composition $\psi \circ \nu$. $\phi$ is not a pseudo-proper map, we will see

\begin{proposition}\label{Prop:Add}\cite[Proposition~4.4.2]{abott}
    $\phi_! \colon \Shv(\Ran) \to\Shv (\Ran_{untl})$
    exists and identifies canonically with $\mathrm{AddUnit}$
\end{proposition}

\begin{remark}
    This functor is the first time to show you the power of the $(\infty,2)-$category.
\end{remark}

\begin{definition}
     $\mathrm{ObvlUnit}\coloneqq \phi^!$
\end{definition}

We need the following lemma to prove Proposition \ref{Prop:Add}.
\begin{lemma}
    \[
    \id_{\Shv(\Ran)} \xrightarrow[]{\sim} (\nu^! \circ \xi^!)
    \]
\end{lemma}
\begin{proof}
    because $\xi \circ \nu \simeq \id$, therefore apply $\Shv$, we get our result.
\end{proof}

\begin{corollary}
    The functor $\nu_!$ identifies to $\xi^!$
\end{corollary}

\begin{proof}
    By the uniqueness of adjunction.
\end{proof}

Therefore we have $\mathrm{AddUnit} = \psi_!\circ \xi^! \simeq \psi_! \circ \nu_! \simeq \phi_!$


\appendix
\printbibliography

\end{document}